
\subsection{The phage--bacteria interaction}
A microbiome is the community of microorganisms living in a particular environment, together with their genetic 
material, interactions, and often the surrounding ecological context. They can be found in soil, in water, and on 
and inside multicellular organisms such as humans. One of the examples of a microbiome is the gut microbiome. 
The gut microbiome contains a dense bacterial community alongside a large population of
viruses that infect some of the bacteria inside the microbiome \citep{shkoporov2019}.

Some of those viruses are bacteriophages (also called phages), meaning viruses that are hosted by bacteria. 
They make up the majority of the gut virome \citep{shkoporov2019}. A phage consists of little
more than a genome enclosed in a protein capsid, in many cases attached to a tail. Infection proceeds in a 
fixed order. First, the proteins at the tip of the tail bind a specific molecular structure on the surface 
of the bacterium, anchoring the phage to the cell. Then, the genome of the phage passes through the tail
into the bacterium while the capsid remains outside. Once inside, the genome is either replicated
with the machinery of the host, which assembles new phages and destroys the cell as they are released, 
or it is inserted into the bacterial chromosome, where it is copied along with the host genome at every cell division until a
later stimulus returns it to the first route. Because the initial binding requires a match between two particular 
molecular structures, a phage can generally infect only a small number of bacterial species, which makes the 
interaction host-specific. 

Phages therefore act continuously on bacteria, removing some members and leaving the rest untouched.
The function of the microbiome as a whole depends on its composition, and since the composition changes with
bacteriophages, its function changes as well. These changes in the composition
of the gut microbiome have been associated with a range of conditions, including colorectal cancer (CRC), 
inflammatory bowel disease (IBD) and graft-versus-host disease (GvHD) \citep{wirbel2019,lloydprice2019,peled2020}.

\subsection{Protein clusters}
Proteins are essential to the interaction between phages and bacteria, as everything from the tail
of the phage that connects to the bacterium to the systems that defend the bacterium against the phage is made
of proteins. The proteins of an organism therefore describe whether it is able to take part in 
an infection.

One of the major problems with using the proteins themselves is that even a single genome carries thousands of proteins,
and that the same protein is rarely exactly the same sequence in two organisms. Two organisms that carry the same kind of 
tail protein would therefore be described by two different entries, and no vocabulary shared between organisms 
would exist. This is solved by grouping protein sequences that are similar enough to each other into protein 
clusters, where every cluster stands for a protein family rather than for one particular variant of it. MMseqs2 \citep{steinegger2017} 
performs this grouping systematically for large collections of sequences. At that size, 
it is not feasible to compare each sequence to all other sequences, so it first searches for short matching subsequences to 
find candidate pairs, scores those cheaply, and aligns only the ones that survive, before forming the clusters 
from the remaining similarities.

After each protein is assigned to a cluster, an organism can be written as a vector that counts how many of
its proteins fall into each cluster. Collecting these vectors gives the feature matrices $X_b$ for the
bacteria and $X_v$ for the viruses, with the organisms as rows and the clusters as columns. Since any single organism
carries only a small share of all clusters that occur across the data, the features are high-dimensional and mostly zero.
Their rows are also genera and vOTUs rather than single genomes, so a profile describes the 
repertoire of a whole taxonomic unit and not that of one organism.

\subsection{CRISPR and SpacePHARER}
CRISPR is an adaptive immune system found in many bacteria. A CRISPR array is a region of the bacterial 
genome built from short repeated sequences that are separated by unique sequences called spacers, together with 
a set of neighbouring \textit{cas} genes \citep{hille2018}. The system operates in three stages. When foreign
DNA enters the cell, proteins encoded by the \textit{cas} genes cut a short fragment out of it and insert that fragment 
into the array as a new spacer. Then, that array is transcribed and processed, so that every spacer ends up in a 
small RNA molecule serving as a guide. If DNA matching one of those guides enters the cell again, it is cut 
and the infection fails. Each spacer is therefore a record of an infection that the bacterial lineage 
encountered in the past, written into its own genome.

The CRISPR mechanism implies that if the spacers of bacteria are compared against the genomes of phages, 
one can determine that the bacterium, or an ancestor of it, was infected, while the absence of a match does 
not establish the opposite. One of the main issues is that spacers are only around 30 nucleotides long, meaning 
they are short enough that a standard nucleotide search loses sensitivity, and a spacer that was acquired long 
ago will have diverged from the phage genome it was taken from. SpacePHARER \citep{zhang2021spacepharer} 
addresses both. It translates the spacers and the phage genomes in all six reading frames and compares them at 
the protein level, where divergence is tolerated better, before realigning the resulting hits at the nucleotide 
level. It then combines the evidence from all spacers of one bacterium against one phage genome into a single 
score, so that several weak matches can be significant together, and reports only the pairs with a false 
discovery rate below 0.05.

There are further caveats to those pairs. Only bacteria that carry a CRISPR system can produce spacers at all, 
and the carriage differs between taxonomic groups. Spacers are also lost over time, so an 
interaction that left no surviving record cannot be distinguished from one that never happened. The record is 
historical rather than current, and spacers are acquired from mobile genetic elements in general and not from 
phages alone. The reported pairs are therefore best read as a lower bound of the interactions that 
happened.

\subsection{Graphical lasso and StARS}
Another signal for phage--bacteria interaction can be obtained from abundances rather than mechanistic evidence.
If the abundances of a bacterium and a phage are measured across a set of samples, and the two organisms interact, 
then their abundances should vary together to some degree. Pairs whose abundances covary
can therefore be treated as candidate edges of a co-abundance network. The problem with using the correlation between two organisms is that
it also reacts to associations that pass through a third one. For example, if two bacteria both increase under the
same condition, and a phage follows one of them, the phage will correlate with the other bacterium
as well, without any relation between the two. Therefore, a correlation network is dense with indirect edges.

The density problem is solved by calculating the correlation between two organisms while holding all the other
organisms constant. This is described by the inverse of the covariance matrix, which is called the precision matrix $\Theta$.
An entry of $\Theta$ that is zero means that the two organisms are conditionally independent given all others, 
so a nonzero entry is a direct association rather than one inherited through a third organism. However, 
estimating $\Theta$ by inverting the sample covariance is not possible when there are more organisms than 
samples. The graphical lasso \citep{friedman2008} solves this by adding an $L_1$ penalty to the Gaussian log-likelihood,

\begin{equation}
\label{eq:glasso}
\hat{\Theta} = \arg\min_{\Theta \succ 0} \left\{ -\log\det\Theta + \mathrm{tr}(\hat{\Sigma}\Theta)
+ \lambda \|\Theta\|_1 \right\} ,
\end{equation}

which forces most entries of $\hat{\Theta}$ to be exactly zero. The penalty weight $\lambda$
controls how many edges remain.

This leaves $\lambda$ to be chosen. Since the network is not predicting anything, there is no prediction error 
to cross-validate, and the usual likelihood-based criteria tend to select too many edges. Instead, 
StARS \citep{liu2010} chooses $\lambda$ via stability. The samples are subsampled
$B$ times, the graphical lasso is fitted on each subsample, and $\lambda$ is selected as the least sparse value at
which the disagreement between the subsampled edge sets stays below a threshold $\beta$. The accepted instability
is controlled by the new threshold, meaning that a larger $\beta$ returns a denser network and a smaller one a
sparser network. As a result, it gives a score for each pair, namely the fraction of the subsamples in which its 
edge was selected,

\begin{equation}
\label{eq:stars}
W_{ij} = \frac{1}{B} \sum_{b=1}^{B} \mathbf{1}\!\left[ \hat{\Theta}^{(b)}_{ij} \neq 0 \right] ,
\end{equation}

which lies between zero and one. Taking only the rows belonging to bacteria and the columns belonging to 
viruses gives a bipartite matrix of the same shape as the CRISPR one.

$W_{ij}$ is not an observation and not a biologically measured quantity,
but the output of an estimator applied to abundance data and a selection frequency rather than a strength of 
association. Two organisms may covary for reasons that have nothing to do with infection, and an edge that is 
selected in most subsamples is evidence that the estimator is stable, not that the two organisms interact.
This caveat is not specific to these data: correlation-based methods recover known infection networks at close
to chance level in simulated virus--microbe communities \citep{coenen2018}, and comparable cautions have been
raised for microbial association networks in general \citep{pinto2022} and for reading viral ecology off
abundance correlations in marine metagenomes \citep{alrasheed2019}.

\subsection{Related work}
Predicting which phage infects which bacterium is an established computational problem, and most existing 
approaches work from genome sequences: they match CRISPR spacers, search for phage sequences inside bacterial 
genomes, or compare the sequence composition of the two organisms, and they are usually evaluated against 
reference collections in which the host of a phage is already known \citep{edwards2016}. A separate line 
of work infers association networks from abundance data, where an edge is a statistical association whose 
ecological meaning is not settled \citep{friedman2008,liu2010}. This thesis sits between the two. It uses 
neither a curated reference collection nor any functional annotation, but describes the organisms by protein 
clusters formed within each cohort and treats CRISPR linkage and co-abundance as two separate prediction targets on the same features, 
within three disease cohorts.

Two recent surveys place this more precisely. Existing methods have been sorted by the features they use,
covering CRISPR spacers, prophage and homology evidence, k-mer composition, viral marker genes and
receptor-binding proteins \citep{nie2024}, and benchmarked against reference collections in which the host of
each virus is already annotated \citep{shang2025}. The closest prior work is RaFAH \citep{coutinho2021}, which 
predicts a host from the protein content of a virus. A second close comparison is iPHoP \citep{roux2023}, which
integrates several such signals for metagenome-derived viruses. Both build their features from curated resources and both return a host taxon
for a virus. The setting here differs in two respects: the protein clusters are formed \emph{de novo} within
each cohort and carry no functional annotation, and a prediction is made for every genus--vOTU pair on a full
grid rather than as a taxon assignment. We are not aware of prior work that combines the two.